\par\noindent\begin{minipage}{\linewidth}
% !TeX root = ../教材题库_全部条目_转写工作稿.tex
% Source: XJ-B1-C03-HW-P086-Y0288
% 第3章 / 3.2 函数的基本性质 / 教材 p.86 / PDF 第93页
% 已对照原页转写题干及公式；待教师复核；未制作答案
% 原题及相邻原解截图（核对用）：../题目原图/XJ-B1-C03-HW-P086-Y0288-01.webp
% 本文件只在转写完成后作为题干使用；OCR 条目尚不是可靠题干。
\begin{bankitem}{习题3.2 第7题}{第3章 / 3.2 函数的基本性质 / 教材 p.86 / PDF 第93页}{XJ-B1-C03-HW-P086-Y0288}{已对照原页转写题干及公式；待教师复核；未制作答案}
已知函数 $f(x)$ 是定义在 $\mathbb R$ 上的偶函数，当 $x\ge0$ 时，$f(x)=2x(x+4)$，试求出函数 $f(x)$ 在 $\mathbb R$ 上的解析式。
\end{bankitem}
\end{minipage}\par\medskip
